
\documentclass{article}
\usepackage{colortbl}
\usepackage{makecell}
\usepackage{multirow}
\usepackage{supertabular}

\begin{document}

\newcounter{utterance}

\centering \large Interaction Transcript for game `codenames', experiment `unambiguous', episode 2 with reflexion{-}Qwen3.5{-}9B.
\vspace{24pt}

{ \footnotesize  \setcounter{utterance}{1}
\setlength{\tabcolsep}{0pt}
\begin{supertabular}{c@{$\;$}|p{.15\linewidth}@{}p{.15\linewidth}p{.15\linewidth}p{.15\linewidth}p{.15\linewidth}p{.15\linewidth}}
    \# & $\;$A & \multicolumn{4}{c}{Game Master} & $\;\:$B\\
    \hline

    \theutterance \stepcounter{utterance}  
    & & & \multicolumn{2}{p{0.3\linewidth}}{
        \cellcolor[rgb]{0.9,0.9,0.9}{
            \makecell[{{p{\linewidth}}}]{
                \texttt{\tiny{[GM$|$GM]}}
                {'hidden': {'team': ['vaudeville', 'peninsula', 'coleslaw', 'disruption'], 'innocent': ['vinegar', 'cocaine', 'propeller', 'bikini'], 'opponent': [], 'assassin': ['census']}, 'revealed': {'team': {'team': [], 'innocent': [], 'opponent': [], 'assassin': []}, 'opponent': {'team': [], 'innocent': [], 'opponent': [], 'assassin': []}}}
            }
        }
    }
    & & \\ \\

    \theutterance \stepcounter{utterance}  
    & & \multicolumn{4}{p{0.6\linewidth}}{
        \cellcolor[rgb]{0.9,0.9,0.9}{
            \makecell[{{p{\linewidth}}}]{
                \texttt{\tiny{[P1$\langle$GM]}}
                \texttt{You have attempted this same task before and did not succeed. These are the plans you wrote after those attempts.} \\
\texttt{Trial \#1: New plan: I failed because I grouped "disruption" with "vaudeville" under the clue "show," which was too ambiguous and allowed the opponent to guess "stead" (a distractor/opponent word) before I could clear "disruption." I also ignored the remaining opponent word "stead" in my final turn, giving them a chance to guess my last target. Next time, I will prioritize clearing the most difficult or isolated words (like "disruption" or "peninsula") with highly specific, unambiguous clues (e.g., "conflict" for disruption) rather than broad categories. I will also explicitly check the remaining opponent word list before sending a clue to ensure I don't inadvertently trigger their guess on a word that could lead to a loss.} \\
\\ 
\texttt{Trial \#2: New plan: I failed because "show" was too broad, linking "vaudeville" and "disruption" in a way that allowed the opponent to guess "stead" (a distractor/opponent word) before I could clear "disruption." I also ignored the remaining opponent word "stead" in my final turn, giving them a chance to guess my last target. Next time, I will prioritize clearing the most difficult or isolated words (like "disruption" or "peninsula") with highly specific, unambiguous clues (e.g., "conflict" for disruption) rather than broad categories. I will also explicitly check the remaining opponent word list before sending a clue to ensure I don't inadvertently trigger their guess on a word that could lead to a loss.} \\
\\ 
\texttt{Trial \#3: New plan: I failed because "conflict" was too broad and only targeted "disruption," leaving "vaudeville," "peninsula," and "coleslaw" vulnerable. The opponent likely used a clue related to "vinegar" (a distractor) or "cocaine" to guess "stead" (an opponent word), which ended my turn prematurely. Next time, I will prioritize clearing the most isolated or high{-}risk words first using highly specific, unambiguous clues (e.g., "sand" for peninsula, "salad" for coleslaw, "comedy" for vaudeville) rather than broad categories. I will also ensure I clear all four of my words before the opponent can make a single move, avoiding any single{-}turn strategy that leaves multiple targets exposed.} \\
\\ 
\\ 
\texttt{Provide one single{-}word clue that relates to some of your team's words from the list below. You can choose to how many and to which words your clue relates to.} \\
\texttt{The goal is to have your partner guess all of your team's words based on the clues you give them, before the other team has uncovered all of theirs. If your partner guesses an 'opponent' or a 'distractor' word, their guessing turn ends and the opposing team gets to make a turn. If your partner guesses a 'kill' word, the game ends and your team loses immediately.} \\
\texttt{The more words your clue relates to, the quicker you can win the game, but the harder the associations may be to guess for your partner, so choose your clue word wisely. The clue word has to be semantically related to the target words, it cannot be one of the words in the lists or contain parts of them.} \\
\\ 
\texttt{Always give your single{-}word clue and your comma{-}separated list of related target words in the following format and make your answers as short as possible, never include any other text than is required in this form:} \\
\\ 
\texttt{CLUE: <WORD>} \\
\texttt{TARGETS: <TARGETS>} \\
\\ 
\texttt{Your team words are: vaudeville, peninsula, coleslaw, disruption.} \\
\texttt{Your opponent's team words are: .} \\
\texttt{Distractor words are: vinegar, cocaine, propeller, bikini.} \\
\texttt{The kill words are: census.} \\
            }
        }
    }
    & \\ \\

    \theutterance \stepcounter{utterance}  
    & \multicolumn{4}{p{0.6\linewidth}}{
        \cellcolor[rgb]{0.9,0.9,0.9}{
            \makecell[{{p{\linewidth}}}]{
                \texttt{\tiny{[P1$\rangle$GM]}}
                \texttt{CLUE: comedy} \\
\texttt{TARGETS: vaudeville} \\
            }
        }
    }
    & & \\ \\

    \theutterance \stepcounter{utterance}  
    & & & \multicolumn{2}{p{0.3\linewidth}}{
        \cellcolor[rgb]{0.9,0.9,0.9}{
            \makecell[{{p{\linewidth}}}]{
                \texttt{\tiny{[GM$|$GM]}}
                \texttt{comedy} \\
            }
        }
    }
    & & \\ \\

    \theutterance \stepcounter{utterance}  
    & & & \multicolumn{2}{p{0.3\linewidth}}{
        \cellcolor[rgb]{0.9,0.9,0.9}{
            \makecell[{{p{\linewidth}}}]{
                \texttt{\tiny{[GM$|$GM]}}
                ['vaudeville']
            }
        }
    }
    & & \\ \\

    \theutterance \stepcounter{utterance}  
    & & & \multicolumn{2}{p{0.3\linewidth}}{
        \cellcolor[rgb]{0.9,0.9,0.9}{
            \makecell[{{p{\linewidth}}}]{
                \texttt{\tiny{[GM$|$GM]}}
                {'word': 'vaudeville', 'assignment': 'team'}
            }
        }
    }
    & & \\ \\

    \theutterance \stepcounter{utterance}  
    & & & \multicolumn{2}{p{0.3\linewidth}}{
        \cellcolor[rgb]{0.9,0.9,0.9}{
            \makecell[{{p{\linewidth}}}]{
                \texttt{\tiny{[GM$|$GM]}}
                \texttt{opponent has won} \\
            }
        }
    }
    & & \\ \\

    \theutterance \stepcounter{utterance}  
    & \multicolumn{4}{p{0.6\linewidth}}{
        \cellcolor[rgb]{0.95,0.95,0.95}{
            \makecell[{{p{\linewidth}}}]{
                \texttt{\tiny{[P2$\rightarrow$P1]}}
                \texttt{Post{-}episode Reflexion (not sent to GM):} \\
\texttt{New plan: I failed because my clue "comedy" was too broad, linking only "vaudeville" while leaving "peninsula" and "coleslaw" exposed, allowing the opponent to guess a distractor or opponent word and end my turn prematurely. Next time, I will prioritize clearing the most isolated or high{-}risk words first using highly specific, unambiguous clues (e.g., "sand" for peninsula, "salad" for coleslaw, "comedy" for vaudeville) rather than broad categories. I will also ensure I clear all four of my words before the opponent can make a single move, avoiding any single{-}turn strategy that leaves multiple targets exposed.} \\
            }
        }
    }
    & & \\ \\

\end{supertabular}
}

\end{document}
